Reasoning demonstrations consume input context every time they are reused. This dissertation studies whether auditable compression can improve the balance between answer accuracy and input use on CLADDER. The procedure, called PNS, freezes and segments a parent trace with a correct final answer and generates candidate continuations from recorded prefix interventions. Shorter demonstrations are selected through answer, semantic-review and length checks. Phase56 accepted shorter traces for 52 of 56 parents, reducing aggregate reasoning tokens across the 52 accepted pairs by 35.74\%. The demonstration bank retains the four original traces for which no candidate was accepted.

Two four-condition evaluations compared PNS with full reasoning, a length-matched heuristic and a zero-shot reference. Demonstration identities, order and count were held fixed across the three conditions using examples. On 246 newly generated numerical instances, PNS answered 224 questions correctly, compared with 208 for each other demonstration condition, while using 18.62\% fewer input tokens than full reasoning. The accuracy gain was 6.50 percentage points, with pointwise 95\% confidence intervals of [2.44, 10.57] percentage points against full reasoning and [2.03, 10.98] against the heuristic. Both gains remained statistically significant after Holm adjustment across six prespecified comparisons. The fresh instances covered four probability-query types and were generated from uniform proposals for conditional probability tables (CPTs), using fixed validity and novelty filters within existing graph families and story templates. The previously evaluated cohort contained 300 targets in 271 groups defined by graph structure and complete CPTs. PNS answered 279 correctly, compared with 277 for full reasoning and 272 for the heuristic, while using 20.19\% fewer input tokens than full reasoning. The CPT-cluster confidence intervals for both accuracy differences included zero, leaving the accuracy differences in this cohort unresolved.

The complete construction and fallback policy improved the accuracy--context trade-off on the fresh cohort and reduced input use in both evaluations. The trace records allowed exact reconstruction, but semantic audits still identified errors in traces with correct final answers. The measured input savings do not establish monetary break-even because construction costs were not fully accounted for.

Keywords. causal reasoning; in-context learning; reasoning compression; trace audit; CLADDER.
